\chapter{Анализ данных}

Ключевой целью данной научно-исследовательской работы является анализ эффективности различных архитектур гибридных силовых установок. В предыдущих разделах была проведена их классификация, которая в дальнейшем используется в качестве основы для сравнительного анализа. Основное внимание уделяется влиянию архитектурных решений трансмиссии и принципов взаимодействия двигателя внутреннего сгорания и электрической машины на топливную экономичность автомобиля.

В рамках данного исследования не рассматриваются модели автомобилей, относящиеся к подключаемым гибридам (PHEV), поскольку их расход топлива в значительной степени определяется ёмкостью аккумуляторной батареи и возможностью внешней подзарядки. Указанные факторы не позволяют корректно выделить вклад самой архитектуры гибридной трансмиссии и могут искажать результаты анализа.

Прямое сравнение расхода топлива различных моделей автомобилей также является некорректным, так как они отличаются по массе, характеристикам двигателя и другим конструктивным параметрам. В связи с этим в работе используется подход, основанный на переходе к безразмерному коэффициенту эффективности, при котором каждая гибридная модель сравнивается с максимально близким по техническим характеристикам аналогом с традиционной силовой установкой. Это позволяет количественно оценить эффект гибридизации при прочих равных условиях и выполнить сопоставление различных архитектур.

В качестве источников данных по расходу топлива были использованы базы ADAC и FuelEconomy.gov (EPA), поскольку они предоставляют стандартизированные и воспроизводимые результаты испытаний, выложенные в открытый доступ на их сайтах. Использование данных источников обеспечивает сопоставимость результатов и минимизирует влияние условий реальной эксплуатации и региональных особенностей испытательных циклов.

\newpage

Ниже приведён список исследуемых гибридов:

\begin{table}[h]
    \centering
    \caption{Архитектуры гибридных автомобилей и их аналоги на ДВС}
    \renewcommand{\arraystretch}{1.25}
    \label{tab:hybrids}
    \begin{tabularx}{\textwidth}{l X l}
        \hline
        Модель гибрида & Модель аналога на ДВС & Источник \\
        \hline
        \textbf{MHEV parallel} \\
        2012 Honda Insight II 1.3 л & 2017 Honda Jazz 1.3 л & ADAC \\
        2024 VW Golf eTSI 1.5 л & 2024 VW Golf TSI 1.5 л & ADAC \\
        2025 Audi A5 TDI 2.0 л & 2017 Audi A5 TFSI 2.0 л & ADAC \\
        \hline
        \textbf{Inline parallel } \\
        2024 BMW 530d H. 3.0 л & 2018 BMW 530d 3.0 л & ADAC \\
        2020 Hyundai Sonata H. 2.0 л & 2019 Hyundai Sonata 2.0 л & EPA \\
        2021 Hyundai Elantra 1.6 л & 2020 Hyundai Elantra 1.6 л & EPA \\
        \hline
        \textbf{HSD} \\
        2022 Toyota Prius 1.8 л & 2022 Toyota Corolla 1.8 л & EPA \\
        2024 Toyota Camry H. 2.5 л & 2024 Toyota Camry LE 2.5 л & EPA \\
        2020 Ford Fusion H.2.0 л & 2020 Ford Fusion 2.0 л & EPA \\
        \hline
        \textbf{MMD series} \\
        2020 Honda Accord H. 2.0 л & 2020 Honda Accord 2.0 л & EPA \\
        2015 Honda Civic H. 1.5 л & 2016 Honda Civic 1.5 л & EPA \\
        2023 Honda Jazz i-MMD 1.5 л & 2018 Honda Jazz i-VTEC 1.5 л & ADAC \\
        \hline
    \end{tabularx}
\end{table}

Как видно, в списке упущены архитектуры REEV и MM-Parallel, так как они встречаются только в конфигурации PHEV. Также отсутствует описанный ранее GM TMH, поскольку данная архитектура оказалась неуспешной и на сегодняшний день не используется в серийных автомобилях.

Все отобранные из источников данные представлены в листинге кода в приложении А.

\newpage

На основе собранных данных по расходу топлива различных моделей гибридных автомобилей, представленных в таблице \ref{tab:hybrids}, отобразим влияние архитектуры гибридной силовой установки на топливную экономичность на рисунке \ref{img:g1}:

\includeimage
    {g1} % Имя файла без расширения (файл должен быть расположен в директории inc/img/)
    {f} % Обтекание (без обтекания)
    {h} % Положение рисунка (см. figure из пакета float)
    {1\textwidth} % Ширина рисунка
    {Обзор влияния архитектуры гибридной силовой установки на топливную экономичность} % Подпись рисунка

Анализ представленных данных показывает, что большинство исследуемых архитектур гибридных силовых установок обеспечивают значительное улучшение топливной экономичности по сравнению с аналогами на традиционных ДВС. Также видно, что в разных условиях движения эффективность гибридных систем варьируется.

Лучшие результаты по снижению расхода топлива демонстрируют архитектуры HSD и MM-Series, которые обеспечивают наиболее значительный прирост топливной экономичности в городском цикле. Это объясняется их способностью эффективно использовать рекуперацию энергии при частых остановках и стартах, характерных для городского движения. Однако на шоссе место первенства уходит Parallel Inline, что связано с прямым потоком мощности от двигателя до колес без двойной конвертации в электрическом контуре. Мягкие гибриды, в свою очередь, показывают умеренное улучшение экономичности, что обусловлено их ограниченными возможностями по электрификации.

\newpage

Выразим средний коэффициент эффективности для каждой архитектуры, что позволит более наглядно сравнить их между собой. Результаты представлены на рисунке \ref{img:g2}:

\includeimage
    {g2} % Имя файла без расширения (файл должен быть расположен в директории inc/img/)
    {f} % Обтекание (без обтекания)
    {h} % Положение рисунка (см. figure из пакета float)
    {0.7\textwidth} % Ширина рисунка
    {Средний рост топливной экономичности} % Подпись рисунка

Наглядно видно преимущество архитектуры HSD перед MM-Series в любых условиях движения. Это связано с высокой степенью интеграции компонентов и оптимизацией работы двигателя внутреннего сгорания в широком диапазоне нагрузок. Также стоит заметить, что на шоссе разрыв между архитектурами сокращается и эффективность падает в целом для всех архитектур. Parallel Inline, в свою очередь, имеет минимальный разрыв по эффективности между городом и шоссе.

\newpage

Полученные результаты предоставляют возможность сравнить эффективность архитектур между собой, выявляя самые экономичные из них, однако стоит посмотреть, как эта экономичность влияет на динамику автомобиля. Зависимость снижения расхода в комбинированном цикле от времени разгона с 60 до 100 км/ч для каждой модели представлена на рисунке \ref{img:g3}:

\includeimage
    {g3} % Имя файла без расширения (файл должен быть расположен в директории inc/img/)
    {f} % Обтекание (без обтекания)
    {h} % Положение рисунка (см. figure из пакета float)
    {1\textwidth} % Ширина рисунка
    {Связь между экономичностью и динамикой} % Подпись рисунка

На основании представленного графика можно сделать несколько выводов. Во-первых, HSD обычно показывает высокий выигрыш в экономии при относительно умеренном разгоне. Архитектуры Inline часто сбалансированы — умеренный выигрыш по расходу, разгон может быть быстрый или средний. MM-Series модели показывают высокую экономию при среднем разгоне, другие меньше экономят, но быстрее ускоряются. Многое зависит от конкретной реализации. MHEV, как правило, обеспечивают быстрый разгон, но снижение расхода умеренное. То есть MHEV помогают динамике, но экономят меньше топлива.